\documentclass[10pt,a4paper]{article}
\usepackage[top=12mm,bottom=12mm,left=14mm,right=14mm,headsep=2mm,footskip=7mm]{geometry}
\usepackage{fontspec}
\usepackage{xeCJK}
\setmainfont{Noto Sans}
\setsansfont{Noto Sans}
\setCJKmainfont{Noto Sans CJK SC}
\setCJKsansfont{Noto Sans CJK SC}
\usepackage{xcolor}
\usepackage{graphicx}
\usepackage{tikz}
\usetikzlibrary{calc,arrows.meta,positioning}
\usepackage{fancyhdr}
\usepackage{hyperref}
\usepackage{ragged2e}
\usepackage{microtype}
\definecolor{P2Paper}{HTML}{FCFAF7}
\definecolor{P2Ink}{HTML}{2D3238}
\definecolor{P2Muted}{HTML}{70757A}
\definecolor{P2BlueLight}{HTML}{D3EAF5}
\definecolor{P2Apricot}{HTML}{F5D9B8}
\definecolor{P2Rose}{HTML}{E8C7D3}
\definecolor{P2Violet}{HTML}{D7D3EA}
\definecolor{P2Blue}{HTML}{5B9FBE}
\definecolor{P2ApricotAccent}{HTML}{D49757}
\definecolor{P2RoseAccent}{HTML}{B97589}
\definecolor{P2VioletAccent}{HTML}{8177A8}
\definecolor{P2Rule}{HTML}{DDE3E7}
\definecolor{P2Panel}{HTML}{FFFDFC}
\pagecolor{P2Paper}
\color{P2Ink}
\setlength{\parindent}{0pt}
\setlength{\parskip}{3.5pt}
\linespread{1.04}
\hypersetup{hidelinks}
\pagestyle{fancy}
\fancyhf{}
\renewcommand{\headrulewidth}{0pt}
\fancyfoot[L]{\fontsize{6.6}{8}\selectfont\color{P2Muted}PROCESS REPORT · WORKFLOW CONSTRUCTION}
\fancyfoot[R]{\fontsize{6.6}{8}\selectfont\color{P2Muted}\thepage}
\newcommand{\eyebrow}[1]{{\fontsize{7.8}{9.4}\selectfont\bfseries\color{P2RoseAccent}#1}\par\vspace{1pt}}
\newcommand{\pagetitle}[1]{{\fontsize{18.5}{21.5}\selectfont\bfseries\color{P2Ink}#1}\par\vspace{3pt}}
\newcommand{\deck}[1]{{\fontsize{9.0}{13.0}\selectfont\color{P2Ink}\RaggedRight #1}\par\vspace{4pt}}
\newcommand{\tracehead}[2]{{\fontsize{7.1}{8.4}\selectfont\bfseries\color{#1}#2}\par}
\newcommand{\tracetitle}[1]{{\fontsize{10.2}{12.2}\selectfont\bfseries\color{P2Ink}#1}\par\vspace{1.5pt}}
\newcommand{\tracebody}[1]{{\fontsize{8.0}{11.2}\selectfont\color{P2Ink}\RaggedRight #1}\par}
\newcommand{\provenance}[1]{\vspace{2pt}{\color{P2Rule}\hrule height .35pt}\vspace{2.5pt}{\fontsize{6.9}{9.2}\selectfont\color{P2Muted}\RaggedRight\textbf{Original Evidence.} #1}\par}
\tikzset{
 userA/.style={rounded corners=1mm,draw=P2ApricotAccent,line width=.55pt,fill=P2Apricot,fill opacity=.34},
 userR/.style={rounded corners=1mm,draw=P2RoseAccent,line width=.55pt,fill=P2Rose,fill opacity=.34},
 gptB/.style={rounded corners=1mm,draw=P2Blue,line width=.55pt,fill=P2BlueLight,fill opacity=.27},
 gptV/.style={rounded corners=1mm,draw=P2VioletAccent,line width=.55pt,fill=P2Violet,fill opacity=.27},
 arrA/.style={-{Stealth[length=2.0mm,width=1.45mm]},draw=P2RoseAccent,line width=.72pt},
 arrB/.style={-{Stealth[length=2.0mm,width=1.45mm]},draw=P2Blue,line width=.72pt},
 arrV/.style={-{Stealth[length=2.0mm,width=1.45mm]},draw=P2VioletAccent,line width=.72pt},
 tagA/.style={circle,minimum size=4.8mm,inner sep=0pt,draw=P2ApricotAccent,line width=.6pt,fill=P2Paper,font=\fontsize{6.7}{7}\selectfont\bfseries},
 tagR/.style={circle,minimum size=4.8mm,inner sep=0pt,draw=P2RoseAccent,line width=.6pt,fill=P2Paper,font=\fontsize{6.7}{7}\selectfont\bfseries},
 tagB/.style={circle,minimum size=4.8mm,inner sep=0pt,draw=P2Blue,line width=.6pt,fill=P2Paper,font=\fontsize{6.7}{7}\selectfont\bfseries}
}

\begin{document}

\eyebrow{WORKFLOW CONSTRUCTION · 20}
\pagetitle{核心 Evidence 不按固定页数压缩：证据决定版式}
\deck{设计 Process Report 时，我很快发现一个问题：有些关键对话天然就比较长，如果为了“控制在两页”而硬裁，反而会把上下文和判断过程裁掉。所以我的要求是先保证Evidence完整，再决定需要几页。}
\begin{tikzpicture}[x=1mm,y=1mm]
  \node[anchor=north west,inner sep=0] (u) at (0,0) {\includegraphics[width=42mm]{assets/user_page_limit_phrase.png}};
  \node[anchor=north west,inner sep=0] (g) at (0,-42) {\includegraphics[width=68mm]{assets/gpt_paging_rule.png}};
  \coordinate (ua) at ($(u.south east)+(-2mm,6mm)$);
  \coordinate (ga) at ($(g.north east)+(-4mm,-14mm)$);
  \draw[arrA] (ua) .. controls (76,-5) and (78,-38) .. (ga);
  \node[tagA] at ($(ua)+(3.5mm,0)$) {1};
  \node[anchor=north west,text width=89mm,align=left] at (88,-2) {
    \tracehead{P2ApricotAccent}{MICRO TRACE 26}\tracetitle{“一次重要交互需要多少页，就使用多少页”}
    \tracebody{我明确反对给核心 Interaction Evidence 预设1--2页上限。GPT随后也把原则改成：一个Decision需要几页就用几页，只要上下文连续、每页都有作用。}
    \vspace{8mm}
    \tracehead{P2Blue}{LAYOUT PRINCIPLE}\tracetitle{Evidence first, layout second}
    \tracebody{所以页数、截图数量和分页都不能先定死。版式应该帮助老师读懂证据，而不是为了版式漂亮去删证据。}
  };
\end{tikzpicture}
\provenance{我的原句与GPT后续分页规则来自同一轮真实Process Report设计讨论；没有改写我的原话。}
\newpage

\eyebrow{WORKFLOW CONSTRUCTION · 21}
\pagetitle{Annotated Evidence Layout：批注位置服从内容，不固定在右栏}
\deck{接着我又发现，如果所有页面都固定成“左边截图、右边批注”，不同长度的对话还是会被同一模板卡住。所以批注位置后来不再固定，可以根据内容放在上、侧或下方。}
\begin{tikzpicture}[x=1mm,y=1mm]
  \node[anchor=north west,inner sep=0] (u) at (0,0) {\includegraphics[width=42mm]{assets/user_general_evidence_phrase.png}};
  \node[anchor=north west,inner sep=0] (g) at (0,-31) {\includegraphics[width=68mm]{assets/gpt_annotation_layout.png}};
  \coordinate (ua) at ($(u.south east)+(-2mm,4mm)$);
  \coordinate (ga) at ($(g.north east)+(-3mm,-20mm)$);
  \draw[arrA] (ua) .. controls (75,-6) and (78,-32) .. (ga);
  \node[tagA] at ($(ua)+(3.5mm,0)$) {1};
  \node[anchor=north west,text width=91mm,align=left] at (86,-2) {
    \tracehead{P2ApricotAccent}{MICRO TRACE 27}\tracetitle{“一般过程也要截图，但不能变成截图流水账”}
    \tracebody{我还补充了一点：一般过程也应该有真实截图，只是没有必要像核心Decision那样每段都展开。GPT因此把核心交互和普通过程分开，普通过程可以用Contact Sheet或简短注释汇总。}
    \vspace{8mm}
    \tracehead{P2VioletAccent}{ANNOTATION GRAMMAR}\tracetitle{批注是解释层，不是固定模板}
    \tracebody{annotation可以根据内容放在上方交代背景、侧边解释变化、下方说明结果。重点是把当前Decision讲清楚，而不是让每页长成完全一样。}
  };
\end{tikzpicture}
\provenance{本页对应我对“一般过程也需要Evidence”的补充，以及GPT随后形成的灵活Annotated Evidence Layout。}
\newpage

\eyebrow{WORKFLOW CONSTRUCTION · 22}
\pagetitle{从“截图”到证据链：Raw、Annotated、Contact Sheet与Final Confirmation}
\deck{截图越来越多以后，我又把“原始证据”和“正式呈现”分开：raw截图完整保存，正式版可以加highlight和箭头；一般过程汇成Contact Sheet；核心Decision还必须保留最后到底确认了什么。}
\begin{tikzpicture}[x=1mm,y=1mm]
  \node[anchor=north west,inner sep=0] (u) at (0,0) {\includegraphics[width=42mm]{assets/user_final_confirm_phrase.png}};
  \node[anchor=north west,inner sep=0] (g1) at (0,-28) {\includegraphics[width=62mm]{assets/gpt_evidence_archive.png}};
  \node[anchor=north west,inner sep=0] (g2) at (66,-28) {\includegraphics[width=54mm]{assets/gpt_contact_sheet.png}};
  \coordinate (ua) at ($(u.south east)+(-2mm,4mm)$);
  \coordinate (ga) at ($(g1.north east)+(-3mm,-17mm)$);
  \draw[arrA] (ua) .. controls (66,-6) and (68,-22) .. (ga);
  \node[tagA] at ($(ua)+(3.5mm,0)$) {1};
  \node[anchor=north west,text width=58mm,align=left] at (119,-2) {
    \tracehead{P2RoseAccent}{DECISION LOCK}\tracetitle{讨论必须有“最终到底确认了什么”}
    \tracebody{核心Decision不能只截“讨论得很激烈”的部分，最后确认也要保留。这样老师才能看到方案是怎样真正收敛下来的。}
    \vspace{7mm}
    \tracehead{P2Blue}{ASSET LAYERS}\tracetitle{Raw负责真实性，Annotated负责可读性}
    \tracebody{raw截图不覆盖、不重画；正式版只在上面叠加highlight、箭头、编号和短注释。一般过程则通过Contact Sheet节省篇幅。}
    \vspace{7mm}
    \tracehead{P2VioletAccent}{ARCHIVE}\tracetitle{正文讲清Decision，Evidence Archive保留完整追溯}
    \tracebody{正文不需要把所有原图都塞进去，但关键结论必须能回到对应的完整Evidence。}
  };
\end{tikzpicture}
\provenance{本页来自“最终确认也必须截图、raw与annotated分离、一般过程Contact Sheet化”的连续真实讨论。}
\end{document}
